\begin{figure}[tbp]
\centering
\begin{tikzpicture}[
  x=1cm,y=1cm,
  every node/.style={font=\small},
  box/.style={draw=linkblue!80!black,rounded corners=2pt,
    text width=6.0cm,align=center,inner sep=8pt},
  input/.style={box,minimum height=16mm},
  rule/.style={box,minimum height=35mm},
  output/.style={box,minimum height=10mm},
  flow/.style={-{Latex[length=2.1mm]},thick,draw=black!70}
]
\node[box,text width=13.3cm,minimum height=12mm] (sample) at (3.7,0)
  {\textbf{裁判}\\[2pt]
   采样 $U,O,\bar\pi,A,T_O,M_O$；令 $S_O=AM_O$。
   只发送下述输入。};

\node[input] (uinput) at (0,-1.9)
  {\textbf{方程解码器}\\
   $(U,A,T_U,S_U)$\\
   $T_U=T_O\bar\pi,\quad S_U=S_O\bar\pi$};
\node[input] (oinput) at (7.4,-1.9)
  {\textbf{投影解码器}\\
   $(O,A,T_O,S_O)$};
\draw[flow] (sample.south -| uinput.north) -- (uinput.north);
\draw[flow] (sample.south -| oinput.north) -- (oinput.north);

\node[rule] (urule) at (0,-4.8)
  {从本解码器的可见建议中选取固定的切片见证\\及其合法答案集\\[2pt]
   $\mathcal A=(z+Z)\cap\{v:e_0^\top v=1\}$。\\[6pt]
   \textit{私有随机性：}\\
   从 $\mathcal A$ 中均匀采样 $a_U$。};
\node[rule] (orule) at (7.4,-4.8)
  {令 $H=\ker A$。选 $M_0$\\使 $AM_0=S_O$。\\[5pt]
   \textit{私有随机性：}\\
   采样均匀的 $\sigma\in K^*$，再以概率\\
   $|\widehat f_\sigma(\Phi)|^2$ 采样 $\Phi$。\\[5pt]
   若 $\sigma|_H\ne0$，取固定的 $h\in H$\\
   使 $\sigma(h)=1$；令 $a_O=\Phi(h)$。};
\draw[flow] (uinput) -- (urule);
\draw[flow] (oinput) -- (orule);

\node[output] (uoutput) at (0,-7.9)
  {$a_U\in L_U,\qquad e_0^\top a_U=1$};
\node[output] (ooutput) at (7.4,-7.9)
  {$a_O\in L_O,\qquad e_0^\top a_O=1$};
\draw[flow] (urule) -- (uoutput);
\draw[flow] (orule) -- (ooutput);

\node[box,text width=8.2cm,minimum height=11mm] (verify) at (3.7,-9.6)
  {\textbf{裁判检验}\quad $\bar\pi a_U=a_O$};
\draw[flow] (uoutput.south) -- (verify.north -| uoutput.south);
\draw[flow] (ooutput.south) -- (verify.north -| ooutput.south);
\end{tikzpicture}
\caption{两个解码器使用独立的随机带。二者都不接收采样到的 $M_O$
或投影 $\bar\pi$；右侧解码器从可见建议中自行选取代表 $M_0$。每个
所示的建议映射都经由齐次坐标包含其仿射截距。转移来的切片只用于
分析这些私有规则：右侧解码器不使用左侧的见证或答案集。若左侧建议
不好，或右侧 $\sigma|_H=0$，该解码器返回任意合法答案。}
\label{fig:advice}
\end{figure}
